\documentclass[a4paper,11pt]{article}

\usepackage[utf8]{inputenc}
\usepackage{enumitem}
\usepackage{hyperref}

%BEGIN LATEX
  \usepackage{core}
  \usepackage{fancyhdr}

  \newcommand{\mousevox}[2]{\IfStrEq{\jobname}{\detokenize{mouse}}{#1}{#2}}
%END LATEX

% me la .xairul. badydi'u
% Hyrule Castle
\title{\vspace{-50pt}\mousevox{
  \profilesvg{mouse} \Huge{
    % Le journal de la souris éclair
    % smacu: x1 is a mouse of species x2
    % ckemono: x1 is an anthropomorphic animal of type/species x2
    % snukarni: k1 is a [shared journal]/blog about topic/subject k2 published/administered by k3 for participants/audience k4.
    % zgavevu: x1 is visible electricity (lightning, electric arcing, etc.) produced by x2
    % te lo zgavevu ckemono smacu ku snukarni
    % te "souris éclair" snukarni

    % lidysmacu: light mouse
    me la lidysmacu pe le snukarni
  }
} {
  \profilesvg{navi} \Huge{
    % Navi's Blog
    % snukarni: k1 is a [shared journal]/blog about topic/subject k2 published/administered by k3 for participants/audience k4.
    me la .navi. pe le snukarni
  }
}}
\author{
  % Navi
  la .navi.
}
\date{2024}

\begin{document}

\maketitle

\thispagestyle{empty}

\subsection*{2024}

\begin{itemize}[itemsep=0\baselineskip,parsep=0pt]
  % Du café façon Guinness
  % Coffee with a tint of dark beer
  %
  % - birje: x1 est de la bière brassée par x2
  % - manku: x1 est sombre
  % - simsa: x1 est semblable à x2 selon x3
  % - ckafi: x1 est du café de type x2
  % 
  %
  % - lo = gadri : Factuel générique « convertit un selbri en sumti ». Le résultat est flou en ce qui concerne l'individualité/groupement
  % - lo'e - gadri : factuel « convertit un selbri en sumti ». Le sumti fait référence à l'archétype de lo {selbri}.
  %
  % - le = gadri : Descriptif générique « convertit un selbri en sumti ». Le résultat est flou en ce qui concerne l'individualité/groupement.
  % - le'e = gadri : Descriptif « convertit un selbri en sumti ». Le sumti fait référence à l'archétype décrit ou perçu de le {selbri}.
  % - poi: débute une clause relative restrictive (ne peut s'attacher qu'à un sumti)
  % - ke'a: sumka'i; se réfère au sumti auquel la clause relative est attachée.
  %
  % help1: lo ckafi poi ke'a simsa lo   manku birje
  % help2: le ckafi poi ke'a simsa le'e manku birje
  %
  % help0: le ckafi poi simsa lo'e manku birje lo ka se tengu makau

  % Aknowledge:
  % https://jboski.lojban.org
  % https://lojban.github.io/ilmentufa/glosser/glosser.htm (see Boxes)
  % https://glosbe.com/en/jbo/dark
  % https://la-lojban.github.io/parser?l=NLPTree&c=compact-all
  \item \bloglink{article00.fr}{Artigo00, le ckafi poi ke'a simsa le'e manku birje}
  \item \bloglink{article01.fr}{Artigo01, Des sous-verres en cuir} % \blogtag{brown}{leather}\blogtag{blue}{laser}\blogtag{black}{ink} % maker
  % matci: x1 is a mat or pad of material x2
  % revyskapi: s1 is a quantity of leather from s2.
  % mi dasni lo cutci be fi loi revyskapi
  \item \bloglink{article02.fr}{Artigo02, De ArchLinux à Nix/NixOS} % #Ops #Nix
\end{itemize}

\end{document}
